% !TeX spellcheck = ru_RU
% !TEX root = vkr.tex

\section*{Заключение}
\label{sec:conclusion}

В ходе работы разработана система для автоматической сегментации и
морфометрических измерений листьев \textit{Sphagnum} в виде настольного
приложения под Windows. Были выполнены следующие задачи:

\begin{enumerate}
    \item Сформулированы требования к разрабатываемой системе.

    \item Спроектирована архитектура и реализовано приложение с тремя режимами
          работы: интерактивная сегментация, параметрические измерения по
          готовым фрагментам и автоматическая сегментация с измерениями.
          Время обработки одного 4K-кадра в автоматическом режиме~---
          $\approx 1{,}0$~с.

    \item Собрана и размечена обучающая выборка для классификатора отбора
          листьев~--- около $3500$ канонизированных фрагментов, в том числе
          $996$ корректных листьев трёх типов \textit{Sphagnum}.

    \item Обучен классификатор отбора для трёх типов листьев
          \textit{Sphagnum}; на классе \texttt{good\_leaf} при гарантированном
          recall $= 0{,}95$ precision составляет $0{,}88\!-\!0{,}92$,
          AP $> 0{,}97$.

    \item Проведено сравнение автоматических измерений с ручной разметкой
          эксперта на $9$ контрольных снимках: MAPE длины = $5{,}5\,\%$,
          MAPE поперечных профилей = $1{,}90\,\%$.
          Прототип апробирован консультантом, качество признано
          удовлетворительным.
\end{enumerate}

\textbf{Артефакты работы:}

\begin{itemize}
    \item Исходный код приложения: \\
          \url{https://github.com/pavelrubanov/segmentation_system}
    \item Размеченный датасет: \\
          \url{https://cloud.mail.ru/public/jvB9/ReNV4TrpS}
    \item Готовая сборка под Windows~64: \\
          \url{https://github.com/pavelrubanov/segmentation_system/releases}
\end{itemize}
